% !PRINTED 219
\weberpage{219}

\begin{center}
\S.\,71.

\gesperrt{Transformation of the cubic equation.}
\end{center}

The \gesperrt{theorem of Hermite} (\S.\,69) suffices, for the cubic equation, to carry out the transformation without further computation.

For this purpose, if we employ homogeneous variables $z_1, z_2$, we have to form simultaneous invariants of the cubic form

\weq{1}{f(z_1,z_2) = a_0z_1^3 + a_1z_1^2z_2 + a_2z_1z_2^2 + a_3z_2^3}

and of the linear form

\weq{2}{T = -t_0z_1 + t_1z_2}

But if we apply the linear substitution

\weq{3}{\begin{gathered} z_1 = \alpha z_1' + \beta z_2' \\ z_2 = \gamma z_1' + \delta z_2' \end{gathered}}

to $T$, there results

\[ -t_0z_1 + t_1z_2 = -t_0'z_1' + t_1'z_2', \]

provided that

\[ \begin{gathered} t_0' = +t_0\alpha - t_1\gamma, \\ t_1' = -t_0\beta + t_1\delta, \end{gathered} \]

or

\weq{4}{\begin{gathered} rt_1 = \alpha t_1' + \beta t_0' \\ rt_0 = \gamma t_1' + \delta t_0'. \end{gathered}}

But this substitution arises from (3) when one replaces $z_1, z_2, z_1', z_2'$ by $rt_1, rt_0, t_1', t_0'$. Now if $I(a,t_1,t_0)$ is a simultaneous invariant of $f$ and $T$, homogeneous and of degree $\nu$ in $t_0, t_1$ and of weight $\lambda$, so that

\[ I(a',t_1',t_0') = r^\lambda I(a,t_1,t_0), \]

then by this permutation it follows that

\[ I(a',z_1',z_2') = r^{\lambda-\nu}I(a,z_1,z_2), \]

i.\,e.\ $I(a,z_1,z_2)$ is a covariant of $f$.

\gesperrt{One therefore obtains all simultaneous invariants of $f$ and $T$ from the covariants of $f$, by replacing $z_1, z_2$ in them by $t_1, t_0$.}

But the covariants of $f$ we have come to know completely in \S.\,62, 63.

After this it is easy to form the cubic equation which, from the substitution \S.\,68, (12),
%CONT>>
